% !TeX program = xelatex
\documentclass[tikz, margin = 5mm]{standalone}

\usepackage{xfrac}
\usepackage[MnSymbol]{mathspec}
\setmainfont[Mapping=tex-text,Ligatures={NoRequired,NoCommon,NoContextual}]{Comic Sans MS}
\setmathfont(Digits,Latin,Greek)[Numbers={Lining,Proportional}]{Comic Sans MS}
\usetikzlibrary{shapes.geometric,decorations,decorations.pathmorphing}
\usetikzlibrary{intersections}
\usetikzlibrary{patterns}
\usetikzlibrary{calc}
\usetikzlibrary{shapes.misc}
\usetikzlibrary{spy}

\pgfdeclarelayer{bg}
\pgfdeclarelayer{bbg}
\pgfsetlayers{bbg, bg, main}

\definecolor{c1}{HTML}{ac0000}
\definecolor{c2}{HTML}{00a0d5}

\newcommand{\abs}[1]{{\left|{#1}\right|}}

\tikzset{
	point/.style = {circle, draw = black, inner sep = 2pt, fill = black},
	empty point/.style = {circle, draw = black, inner sep = 2pt, fill = white},
	cross/.style = {cross out, draw, minimum size = 2 * (#1 - \pgflinewidth), inner sep = 0pt, outer sep = 0pt},
	xkcd/.style = {decorate, decoration = {random steps, pre length = 4mm, post length = 4mm, segment length = 0.6mm, amplitude = 0.2pt}}}

\begin{document}
	\begin{tikzpicture}
		\draw[xkcd, thick, ->] (-0.5,0) -- (8,0) node[pos=1, below]{$x$};
		\draw[xkcd, thick, ->] (0,-0.5) -- (0,5) node[pos=1, left]{$y$};
		\draw[xkcd, c1, thick, samples = 200, name path global = f] plot [smooth] coordinates {(-0.5,0.8) (1,2.8) (2.1,1.5) (3.6,4.3) (5.5,2.2) (8,0.7)} node[right] {$C_f$};
		\draw[draw = none, name path global = line a] (1,0) -- (1,5);
		\draw[draw = none, name path global = line b] (6.2,0) -- (6.2,5);
		\node[draw = none] (o) at (0,0) {};
		\node[draw = none] (ur) at (8,5) {};
		\node[point, c1, name intersections = {of = f and {line a}}] (a) at (intersection-1) {};
		\node[point, c1, name intersections = {of = f and {line b}}] (b) at (intersection-1) {};
		\draw[xkcd, dashed, thick, c1] (a) -- (a |- o) node [pos=1, below] {$x=a$};
		\draw[xkcd, dashed, thick, c1] (b) -- (b |- o) node [pos=1, below] {$x=b$};
		\begin{scope}
			\clip (1,-0.5) rectangle (b |- ur);
			\fill[xkcd, c1, fill opacity = 0.1, samples = 200] plot [smooth] coordinates {(-0.5,0.8) (1,2.8) (2.1,1.5) (3.6,4.3) (5.5,2.2) (8,0.7)} -- (7.6,0) -- (-0.44,0) -- cycle;
		\end{scope}
	\end{tikzpicture}
\end{document}